\documentclass[11pt]{article}
\usepackage[margin=1in]{geometry}
\usepackage{amsmath,amssymb,amsthm,mathtools}
\usepackage{microtype}
\usepackage{booktabs,longtable,array}
\usepackage[hidelinks]{hyperref}
\usepackage{xcolor}
\usepackage{enumitem}
\usepackage{fancyhdr}
\usepackage{listings}
\usepackage[T1]{fontenc}
\usepackage{lmodern}

\newtheorem{theorem}{Theorem}[section]
\newtheorem{lemma}[theorem]{Lemma}
\newtheorem{proposition}[theorem]{Proposition}
\newtheorem{corollary}[theorem]{Corollary}
\newtheorem{definition}[theorem]{Definition}
\newtheorem{remark}[theorem]{Remark}
\newcommand{\1}{\mathbf 1}
\newcommand{\A}{\mathcal A}
\newcommand{\D}{\mathcal D}
\newcommand{\C}{\mathcal C}
\newcommand{\Ree}{\operatorname{Re}}
\newcommand{\RH}{\mathrm{RH}}
\newcommand{\logp}{\log_+}

\pagestyle{fancy}
\setlength{\headheight}{14pt}
\fancyhf{}
\lhead{T--99000: C4MBI67 primitive-prefix reconstruction}
\rhead{\thepage}

\lstset{basicstyle=\ttfamily\small,breaklines=true,frame=single,columns=fullflexible}

\title{\textbf{Primitive-Prefix Reconstruction of the T--97701/C4MBI67 Route}\\
\large A self-contained exact reduction, a billion-endpoint rational certificate,\\
and the remaining all-scale closure theorem}
\author{Agentic Polymath Project --- Riemann route}
\date{19 August 2026}

\begin{document}
\maketitle

\begin{center}
\fbox{\parbox{0.92\textwidth}{
\textbf{Scientific status.}  Every finite identity, the Mellin--Landau
consumer, and the directed certificate through $10^9$ are proved below.
The eventual all-scale sign of $\A_X$ is \emph{not} proved.  Consequently this
manuscript is not a proof of the Riemann Hypothesis.  It reconstructs the
entire conditional chain without hidden imports and isolates one exact
producer.}}
\end{center}

\begin{abstract}
The T--97701 packet reduces the factor--67 root Bellman inequality to the sign
of an annular $5{:}3$ scalar.  We reconstruct that reduction from the source
dictionary, rather than importing it.  The central new object is the primitive
prefix
\[
 C(t)=6-\left(6B_{1/2}(t)-\frac9{\sqrt2}B_{1/2}(t/2)
                    +\frac32B_{1/2}(t/4)\right),
\]
for which the desired scalar is the logarithmic window
$\A_X=\int_{X/4}^{X}C(t)\,dt/t$.  We give three exact descriptions of $C$:
a Möbius prefix, a generalized-prime inclusion--exclusion sum, and the
expected Euler characteristic of a random weighted threshold complex.  These
yield an exact windowed Bellman recurrence that displays the surviving
correlated-overshoot obstruction.

A linear least-prime-factor sieve combined with exact $2^{-40}$ rational
brackets for every $1/\sqrt n$ proves $C(t)\ge0$ for every real
$1\le t<10^9+1$.  The global lower enclosure is
$1226685126915/2^{40}>1.11565$ and occurs at $t=48433$.  Finally, we prove from
scratch the specialized Landau theorem and the zeta functional-equation facts
needed to show that eventual nonnegativity of $\A_X$ implies RH.  The only
unproved conclusion-producing statement is the eventual all-scale sign itself.
\end{abstract}

\tableofcontents

\section{The exact theorem ledger}

\begin{longtable}{@{}p{0.53\textwidth}p{0.39\textwidth}@{}}
\toprule
Statement & Status\\
\midrule
Source dictionary and Möbius inverse & proved exactly\\
Primitive-prefix/window identity & proved exactly\\
Four-band C4MBI67 identity & proved exactly\\
Logarithmic prime ownership & proved exactly\\
Weighted threshold-complex representation & proved exactly\\
Windowed Bellman recurrence & proved exactly\\
Mellin transform and off-line-zero noncancellation & proved exactly\\
Specialized Landau theorem & proved below\\
Theta-functional-equation symmetry & proved below\\
$C(t)\ge0$ for $1\le t<10^9+1$ & exact rational certificate\\
$\A_X\ge0$ for all sufficiently large $X$ & \textbf{open / RH-bearing}\\
Riemann Hypothesis & \textbf{unproved}\\
\bottomrule
\end{longtable}

\section{Arithmetic reconstruction from the source dictionary}

\subsection{Möbius inversion}

The Möbius function is defined by $\mu(1)=1$, by $\mu(n)=0$ if a prime square
divides $n$, and by $\mu(p_1\cdots p_k)=(-1)^k$ for distinct primes.

\begin{lemma}[Dirichlet inverse of the constant-one function]
For every positive integer $n$,
\[
 \sum_{d\mid n}\mu(d)=\1_{n=1}.
\]
\end{lemma}
\begin{proof}
If $n=1$ the sum is $1$.  If $n>1$ has distinct prime divisors
$p_1,\ldots,p_r$, only squarefree divisors contribute and the sum is
\[
 \sum_{S\subseteq\{1,\ldots,r\}}(-1)^{|S|}=(1-1)^r=0.
\]
\end{proof}

Define the positive sharp source dictionary
\begin{equation}\label{eq:qdict}
q(1)=0,\qquad q(2)=15,\qquad q(4)=3,\qquad
q(n)=6\quad(n\notin\{1,2,4\}).
\end{equation}
Let $a=q*\mu$ under Dirichlet convolution.

\begin{lemma}[Explicit native coefficients]\label{lem:aexplicit}
For every $n\ge1$,
\begin{equation}\label{eq:aexplicit}
 a(n)=6\1_{n=1}-6\mu(n)
 +9\1_{2\mid n}\mu(n/2)
 -3\1_{4\mid n}\mu(n/4).
\end{equation}
\end{lemma}
\begin{proof}
Write
\[
 q=6\mathbf 1-6\delta_1+9\delta_2-3\delta_4,
\]
where $\mathbf1(n)=1$ and $\delta_j(n)=\1_{n=j}$.  Convolution with $\mu$,
using $\mathbf1*\mu=\delta_1$, gives
\[
 q*\mu=6\delta_1-6\mu+9(\delta_2*\mu)-3(\delta_4*\mu),
\]
and $(\delta_j*\mu)(n)=\1_{j\mid n}\mu(n/j)$.
\end{proof}

For later checking, if $n=2^e m$ with $m$ odd and squarefree, then except for
$a(1)=0$,
\begin{equation}\label{eq:aclass}
 a(m)=-6\mu(m),\quad a(2m)=15\mu(m),\quad
 a(4m)=-12\mu(m),\quad a(8m)=3\mu(m),
\end{equation}
and $a(2^e m)=0$ for $e\ge4$ or nonsquarefree $m$.

